% Part: many-valued-logic
% Chapter: syntax-and-semantics
% Section: semantic-notions

\documentclass[../../../include/open-logic-section]{subfiles}

\begin{document}

\olfileid{mvl}{syn}{sem}

\olsection{चिन्हार्थविषयक संकल्पना}

बहुमूल्य तर्कशास्त्र~$\Log L$ एखाद्या मॅट्रिक्सने दिलेले
आहे असे समजू. मग~$\Log L$ साठी नेहमीच्या चिन्हार्थविषयक
संकल्पना परिभाषित करता येतात.

\begin{defn}
\begin{enumerate}
\item !!^a{formula}~$!A$ \emph{पूर्ततायोग्य} असते, जर
  एखाद्या~$\pAssign{v}$ साठी $\pSat{v}{!A}$;
  कोणत्याही $\pAssign{v}$ साठी $\pSat{v}{!A}$ नसेल, तर ते
  \emph{अपूर्ततायोग्य} असते;
\item !!^a{formula}~$!A$ हे \emph{सर्वतः सत्य सूत्र}
  असते, जर सर्व !!{valuation}s~$v$ साठी $\pSat{v}{!A}$;
\item $\Gamma$ हा !!{formula}s चा संच असेल, तर
  $\Gamma \Entails !A$ (``$\Gamma$ पासून $!A$
  निष्पन्न होते'') जर आणि फक्त जर
  $\pSat{v}{\Gamma}$ असलेल्या प्रत्येक
  !!{valuation}~$\pAssign{v}$ साठी $\pSat{v}{!A}$;
\item $\Gamma$ हा !!{formula}s चा संच असेल, तर
  $\pSat{v}{\Gamma}$ असलेले !!a{valuation}~$\pAssign{v}$
  अस्तित्वात असल्यास $\Gamma$ \emph{पूर्ततायोग्य} असतो;
  अन्यथा $\Gamma$ \emph{अपूर्ततायोग्य} असतो.
\end{enumerate}
\end{defn}

अभिजात तर्कशास्त्रात जशी काही तथ्ये या संकल्पनांसाठी
लागू होतात, तशीच ती येथेही लागू होतात:

\begin{prop}
\ollabel{prop:semanticalfacts}
\begin{enumerate}
\item $!A$ हे सर्वतः सत्य सूत्र असते, जर आणि फक्त जर
  $\emptyset \Entails !A$;
\item $\Gamma$ पूर्ततायोग्य असेल, तर $\Gamma$ चा प्रत्येक सांत
  उपसंचही पूर्ततायोग्य असतो;
\item\ollabel{def:monotonicity}%
एकस्वनिकता: $\Gamma \subseteq \Delta$ आणि
  $\Gamma \Entails !A$ असल्यास $\Delta \Entails !A$ देखील;
\item\ollabel{def:Cut}%
संक्रमकता: $\Gamma \Entails !A$ आणि
  $\Delta \cup \{ !A\} \Entails !B$ असल्यास
  $\Gamma \cup \Delta \Entails !B$;
\end{enumerate}
\end{prop}

\begin{proof}
सराव.
\end{proof}

\begin{prob}
\olref[mvl][syn][sem]{prop:semanticalfacts} सिद्ध करा.
\end{prob}

अभिजात तर्कशास्त्रात निष्पन्नता आणि सशर्त संयोजक यांचा संबंध
जोडता येतो. उदाहरणार्थ, \emph{विध्यात्मक अनुमान} वैध असते:
जर $\Gamma \Entails !A$ आणि $\Gamma \Entails !A \lif !B$
असेल, तर $\Gamma \Entails !B$. अभिजात तर्कशास्त्रात
$\Entails$ आणि $\lif$ यांच्यातील दुसरा महत्त्वाचा संबंध म्हणजे
चिन्हार्थविषयक निगमन प्रमेय: $\Gamma \Entails !A \lif !B$
जर आणि फक्त जर $\Gamma \cup \{!A\} \Entails !B$.
हे निष्कर्ष बहुमूल्य तर्कशास्त्रांत \emph{नेहमीच लागू होतात असे नाही}.
ते लागू होणे $\tf{\lif}$ या सत्यता-फलनावर अवलंबून असते.

\end{document}
